\documentclass{article}
%packages
\usepackage{fullpage}
\usepackage{xspace}
\usepackage{graphicx,color}
\usepackage{hyperref}
\usepackage{graphicx}
\usepackage[dvipsnames]{xcolor}
\usepackage{amsmath, amsthm, amscd, amsfonts, amssymb}
\usepackage{graphicx}
\usepackage{stmaryrd}
\usepackage{quiver}
\usepackage{float}
\usepackage{mathtools}
\usepackage{mathrsfs}
\usepackage{multicol}
\usepackage{tikz-cd}
\usepackage[shortlabels]{enumitem}
\usepackage{stackrel}
\usepackage{cancel}
\usepackage[framemethod = TikZ]{mdframed}

%theorems
\newtheorem{theorem}{Teorema}[subsection]

\newtheorem{algorithm}{Algorithm}
\newtheorem{axiom}{Axioma}
\newtheorem{case}{Case}
\newtheorem{claim}{Afirmación}
\newtheorem{conclusion}{Conclusión}
\newtheorem{condition}{Condition}
\newtheorem{conjecture}[theorem]{Conjetura}
\newtheorem{corollary}[theorem]{Corolario}
\newtheorem{criterion}{Criterion}
\newtheorem{proposition}[theorem]{Proposición}
\newtheorem{lemma}[theorem]{Lema}

\theoremstyle{definition}
\newtheorem{definition}[theorem]{Definición}

\newtheorem{exercise}{Ejercicio}

\newtheorem*{notation}{Notación}
\newtheorem{problem}{Problema}[section]
\newtheorem*{obs}{Observación}
\newtheorem*{addendum}{Addendum}

\newtheorem*{tarea}{\wea{Tarea}}

\theoremstyle{remark}
\newtheorem{remark}[theorem]{Remark}
\newtheorem*{solution}{\textbf{Solución}}
\newtheorem{summary}{Summary}
\newtheorem{example}[theorem]{\textbf{Ejemplo}}

\numberwithin{equation}{section}

\renewcommand*{\proofname}{\textbf{Demostración}}

\surroundwithmdframed[
  topline=false,
  rightline=false,
  bottomline=false,
  leftmargin=\parindent,
  skipabove=\medskipamount,
  skipbelow=\medskipamount
]{proof}

\surroundwithmdframed[
  topline=false,
  rightline=false,
  bottomline=false,
  leftmargin=\parindent,
  skipabove=\medskipamount,
  skipbelow=\medskipamount
]{solution}

%commands
\newcommand{\A}{\mathbb{A}}
\newcommand{\B}{\mathcal{B}}
\newcommand{\C}{\mathbb{C}}
\newcommand{\D}{\mathbb{D}}
\newcommand{\F}{\mathbb{F}}
\newcommand{\N}{\mathbb{N}}
\renewcommand{\O}{\mathcal{O}}
\renewcommand{\P}{\mathcal{P}}
\newcommand{\Q}{\mathbb{Q}}
\newcommand{\R}{\mathbb{R}}
\renewcommand{\S}{\mathbb{S}}
\newcommand{\U}{\mathcal{U}}
\newcommand{\Z}{\mathbb{Z}}

\newcommand{\Acal}{\mathcal{A}}
\newcommand{\Bcal}{\mathcal{B}}
\newcommand{\Ccal}{\mathcal{C}}
\newcommand{\Dcal}{\mathcal{D}}
\newcommand{\Ecal}{\mathcal{E}}
\newcommand{\Fcal}{\mathcal{F}}
\newcommand{\Gcal}{\mathcal{G}}
\newcommand{\Hcal}{\mathcal{H}}
\newcommand{\Ical}{\mathcal{I}}
\newcommand{\Jcal}{\mathcal{J}}
\newcommand{\Kcal}{\mathcal{K}}
\newcommand{\Lcal}{\mathcal{L}}
\newcommand{\Mcal}{\mathcal{M}}
\newcommand{\Ncal}{\mathcal{N}}
\newcommand{\Qcal}{\mathcal{Q}}
\newcommand{\Rcal}{\mathcal{R}}
\newcommand{\Scal}{\mathcal{S}}
\newcommand{\Tcal}{\mathcal{T}}
\newcommand{\Ucal}{\mathcal{U}}
\newcommand{\Vcal}{\mathcal{V}}
\newcommand{\Wcal}{\mathcal{W}}
\newcommand{\Xcal}{\mathcal{X}}
\newcommand{\Ycal}{\mathcal{Y}}
\newcommand{\Zcal}{\mathcal{Z}}

\renewcommand{\a}{\alpha}
\renewcommand{\b}{\beta}
\renewcommand{\d}{\delta}
\newcommand{\e}{\varepsilon}
\renewcommand{\t}{\tau}
\newcommand{\lam}{\lambda}
\newcommand{\Lam}{\Lambda}

\newcommand{\vphi}{\varphi}
\newcommand{\oo}{\infty}
\newcommand{\isom}{\cong}
\newcommand{\longto}{\longrightarrow}
\newcommand{\embbed}{\hookrightarrow}

\renewcommand{\bar}{\overline}
\newcommand{\sgn}{\mathrm{sgn}}
\newcommand{\id}{\mathrm{id}}
\newcommand{\im}{\mathrm{im}\,}
\newcommand{\dist}{\mathrm{dist}}

\newcommand{\normal}{\unlhd}

\renewcommand{\tilde}{\widetilde}
\renewcommand{\hat}{\widehat}

\newcommand{\abs}[1]{\left| #1 \right|}
\newcommand{\norm}[1]{\left\lVert #1 \right\rVert}
\newcommand{\br}[1]{\left( #1 \right)}
\newcommand{\set}[1]{\left\{ #1 \right\}}
\newcommand{\cor}[1]{\left[ #1 \right]}
\newcommand{\gen}[1]{\left\langle #1 \right\rangle}
\newcommand{\matr}[1]{\begin{pmatrix}#1\end{pmatrix}}

\newcommand{\znz}[1]{\frac{\Z}{#1 \Z}}
\newcommand{\zmz}[1]{\Z/ #1 \Z}

% Por si me equivoco al escribir
\newcommand{\rac}[2]{\frac{#1}{#2}}
\newcommand{\olon}{\colon}
\newcommand{\ubseteq}{\subseteq}
\newcommand{\ubset}{\subset}
\newcommand{\irc}{\circ}
\renewcommand{\o}{\to}

\newcommand{\ds}{\displaystyle}

\renewcommand*\contentsname{Contenidos}

\newcommand{\wea}[1]{\textcolor{BlueViolet}{\textbf{\emph{#1}}}}
\newcommand{\fecha}[1]{\begin{flushright} \underline{#1} \end{flushright}}

\newcommand{\dem}[1]{\textbf{Demostración (#1)}}

%title
\title{MPG3100 -- Análisis I}
\author{Felipe Inostroza \\ Prof. Mariel Sáez}
\date{Primer Semestre 2026}

\begin{document}

\maketitle
\tableofcontents

\newpage
\fecha{Clase 1. Lunes 9 de Marzo}
\section{Topología}
\subsection{Topologías y espacios métricos}

\begin{definition}
  Dado un conjunto $X$, un conjunto
  \[\tau \subseteq \P(X) = \{\O \colon \O \subseteq X\}\]
  se dice \wea{topología en $X$} si
  \begin{enumerate}[(1)]
    \item $\varnothing, X \in \tau$;
    \item Si $\O_1, \O_2 \in \tau \implies \O_1 \cap \O_2 \in \tau$;
    \item $\{\O_\a\}_{\a\in I} \subseteq \tau \implies \bigcup_{\a \in I} \O_\a \in \tau$.
  \end{enumerate}
\end{definition}

\begin{notation}\
  \begin{itemize}
    \item Si $\O \in \t$, $\O$ se llama \wea{abierto}.
    \item $(X,\tau)$ es un \wea{espacio topológico}.
    \item Si $a \in X$ y $\O \in \t$ tal que $a \in \O$, $\O$ se llama \wea{vecindad de $a$}.
  \end{itemize}
\end{notation}

\begin{example}\
  \begin{enumerate}[(1)]
    \item Para cualquier conjunto $X$ se tiene que $\t = \{\varnothing,X\}$ es una topología de $X$ (\wea{topología trivial}).
    \item $\t = \P(X)$ es la \wea{topología discreta}.
    \item Si $(X,d)$ es un espacio métrico, una topología es
    \[\t = \{\O \subseteq X \colon \forall a \in \O, \exists r > 0 \colon B_r(A) \subseteq \O\}.\]
    Esto incluye $\R^n$, espacios normados, espacios de Hilbert, etc.
  \end{enumerate}
\end{example}

\begin{proposition}
  $E \ubseteq X$ es abierto si y solo si $\forall a \in E$, $\exists \O_a$ vecindad de $a$ tal que $\O_a \subseteq E$.
\end{proposition}
\begin{proof}\
  \begin{itemize}
    \item \boxed\implies
    Si $a \in E$ y $E$ es abierto, entonces $E$ es vecindad de $a$ y $E \subseteq E$. \checkmark
    \item \boxed\impliedby
    Por demostrar: $E$ abierto.
    Si $a \in E$, por hipótesis $\exists \O_a$ tal que $a \in \O_a \subseteq E$, y por lo tanto
    \[\bigcup_{a \in E} \O_a \subseteq E \subseteq \bigcup_{a \in E} \O_a.\]
    Por lo tanto $E = \bigcup_{a \in E} \O_a$, y al ser unión de abiertos, tenemos que $E$ es abierto.
  \end{itemize}
\end{proof}

\begin{definition}
  La \wea{clausura} de un conjunto $E$ está dada por
  \[\bar{E} = \{a \in X \colon \forall \O_a \ni a, \ \O_a \cap E \neq \varnothing\}.\]
\end{definition}

\begin{definition}
  $E$ es cerrado si y solo si $\bar{E} = E$.
\end{definition}

\begin{proposition}
  $E$ es cerrado si y solo si $X \setminus E$ es abierto.
\end{proposition}
\begin{proof}
  \
  \begin{itemize}
    \item \boxed{\implies} Sea $a \in X \setminus E$. Si $\forall \O_a$ vecindad de $a$
    \begin{align*}
      \O_a \not\subseteq X \setminus E
      &\iff \exists b \in \O_a \cap E \\
      &\iff \O_a \cap E \neq \varnothing \\
      &\iff a \in \bar E = E.
    \end{align*}
    Pero esto contradice que $a \in X \setminus E$. Entonces existe $\O_a$ tal que $\O_a \subseteq X \setminus E$, por lo que $X \setminus E$ es abierto.

    \item \boxed{\impliedby} Por demostrar: $\bar E = E$.
    Si $a \in E$, entonces $\forall \O_a$ vecindad de $a$ se tiene que $$a \in \O_a \cap E$$ y entonces $E \subseteq \bar E$. Veamos ahora que $\bar E \subseteq E$.

    Sea $a \in \bar E$, entonces para toda vecindad $\O_a$ se tiene que
    \[\O_a \cap E \neq \varnothing. \tag{$\ast$}\]
    Supongamos que $a \in \bar E \setminus E$, entonces $a \in X \setminus E$. Como $X \setminus E$ es abierto, existe $\O_*$ vecindad de $a$ tal que $\O_* \subseteq X \setminus E$. Entonces $\O_* \cap E = \varnothing$, lo que contradice $(*)$.
    Concluimos que $\bar E \setminus E = \varnothing$.
  \end{itemize}
\end{proof}

\begin{definition}
  [Topología inducida]
  Si $(X,\t)$ es un espacio topológico e $Y \subseteq X$, la \wea{topología inducida en $Y$} es
  \[\t|_Y = \{\O \cap Y \colon \O \in \t\}.\]
\end{definition}

\begin{exercise}
  Verificar que $\t|_Y$ es topología.
\end{exercise}

\begin{obs}
  Con la topología inducida, los abiertos de $Y$ no son necesariamente abiertos en $X$.
\end{obs}

\begin{example}
  En $\R$ con la topología usual, si tomamos $Y = [0,2]$ entonces $[0,1)$ es abierto, ya que $[0,1) = [0,2] \cap (-1,1)$ y $(-1,1)$ es abierto en $\R$.
\end{example}

\begin{example}
  Consideremos la \wea{2--esfera}
  \[\S^2 = \{(x,y,z) \in \R^3 \colon x^+ y^2 + z^= 1\}.\]
  El conjunto
  \[\{(x,y,z) \in \S^2 \colon z>0\}\]
  es abierto en la topología inducida por $\R^3$.
\end{example}

\begin{definition}
  Sea $P \in X$. Decimos que $\B_p \subseteq \t$ es una \wea{base de la topología $\t$ en $p$} si
  \begin{enumerate}[(i)]
    \item Todos los elementos de $\B_p$ son vecindades de $p$;
    \item Para toda vecindad $U \ni p$, existe $B \in \B_p$ tal que $B \subseteq U$.
  \end{enumerate}
  Decimos que $\B$ es \wea{base de la topología} si $\forall p \in X$, $\exists \B_p \subseteq \B$ base de la topología en $p$.
\end{definition}

\begin{example}
  Si $(X,d)$ es un espacio métrico, entonces
  \[\{B_r(p) \colon r \in (0,\oo)\}\]
  es base en $p$.
  \[\{B_r(p) \colon r \in (0,\oo), p \in X\}\]
  es base de $\t$.
\end{example}

\begin{example}
  $X$ con la topología discreta,
  \[\{\{p\} \colon p \in X\}\]
  es base.
\end{example}

\begin{proposition}
  Sea $X \neq \varnothing$ y $\B \subseteq \P(X)$.
  $\B$ es base de alguna topología si y solo si
  \begin{enumerate}[(i)]
    \item $\B$ cubre $X$
    \item $B_1,B_2 \in \B $, $\forall p \in B_1 \cap B_2, \exists B_p \in \B$ tal que $p \in B_p \subseteq B_1 \cap B_2$.
  \end{enumerate}
\end{proposition}
\begin{proof}
  \
  \begin{itemize}
    \item \boxed{\implies}
    $(X,\t)$ es topología y $\B$ es base de $\t$.
    Como $\B$ es base para $p \in X$, entonces existe $B_p \in \B$ vecindad de $p$. Luego
    \[X \subseteq \bigcup_{p \in X} B_p \subseteq X\]
    por lo que $(i)$ se cumple, ya que
    \[X = \bigcup_{p \in X} B_p \subseteq \bigcup_{B \in \B} B \subseteq X.\]

    Veamos ahora (ii). Tomamos $B_1,B_2 \in \B$, entonces $B_1 \cap B_2$ es abierto:
    Si $B_1 \cap B_2 = \varnothing$, entonces ta está en la topología.
    Si $B_1 \cap B_2 \neq \varnothing$, entonces existe $p \in B_1 \cap B_2$. Por definición de base, existe $B$ vecindad de $p$ con $B \subseteq B_1 \cap B_2$.

    \item \boxed{\impliedby}
    Sea $\t = \{\varnothing\} \cup \{\bigcup_{\a \in I} B_\a \colon \{B_\a\}_{\a \in I} \subseteq \B\}$ el conjunto de uniones arbitrarias de elementos de $\B$. Veamos que $\t$ es topología y $\B$ es base de $\t$.
    \begin{itemize}
      \item $\varnothing \in \t$.
      \item Por (i), $X = \bigcup_{B \in \B} B \in \t$, y entonces $X \in \t$.
      \item Si tomamos $\O_\a = \bigcup_{\beta \in I_\a} B_\beta^\a \in \t$
      \[\bigcup_{\beta \in J} \O_\beta = \bigcup_{\a \in I} \bigcup_{\beta \in I_\a} B_\beta^\a \in \t.\]
      \item Sean $\O_1, \O_2 \in \t$, entonces
      \[\O_1 = \bigcup_{\a \in I_1} B_\a^1, \quad
      \O_2 = \bigcup_{\beta \in I_2} B_\b^2.\]
      Sea $p \in \O_1 \cap \O_2$, entonces $p \in \O_1$ y $p \in \O_2$.
      Luego existe $\a \in I_1$ tal que $p \in B_\a^1$ y existe $\b \in I_2$ tal que $p \in B_\b^2$. Por lo tanto $p \in B_\a^1 \cap B_\b^2$.
      Luego $\O_1 \cap \O_2 \in \t$.
    \end{itemize}
  \end{itemize}
\end{proof}

\begin{example}
  Si $(X,\t_X)$ e $(Y,\t_Y)$ son espacios topológicos, entonces
  \[\{\O_1 \times \O_2 \colon \O_1 \subseteq \t_X, \O_2 \subseteq \t_Y\}\]
  es una base para una topología en $X \times Y$.
\end{example}

\begin{definition}
  La topología dada por la base anterior se llama \wea{topología producto}.
\end{definition}

\begin{example}
  La topología producto en $\R \times \R$ es igual a la topología en $\R^2$ dada por la norma usual.
\end{example}

\fecha{Clase 2. Miércoles 11 de Marzo}
\begin{exercise}
  Sea $(X,d)$ espacio métrico. Muestre que $x \in \bar{E}$ si y solo si $\exists (x_n)_{n\in\N} \subseteq E$ tal que
  \[\lim_{n\to\oo} x_n = x.\]
\end{exercise}

\begin{definition}
  Sea $(X,\t)$ un espacio topológico. Si $C \subseteq X$, decimos que $\O$ es \wea{vecindad de $C$} si $\O\in\t$ y $C \subseteq \O$.
\end{definition}

\begin{definition}
  Sea $(X,\t)$ espacio topológico y $(a_n)_{n\in\N} \subseteq X$. Decimos que $\lim_{n\to\oo} a_n = a$ si $\forall \O$ vecindad de $a$, existe $n_0$ tal que $a_n \in \O$, $\forall n \geq n_0$.
\end{definition}

\begin{example}
  Si consideramos $X$ con la topología discreta, entonces una sucesión $(a_n)_{n\in\N}$ es convergente si y solo si $\exists n_0$ tal que $a_n = a$ para todo $n \geq n_0$. Es decir, las sucesiones que convergen son constantes desde un cierto punto en adelante.
\end{example}

\newpage
\subsection{Axiomas de Separación}

\begin{definition}
  Sea $(X,\t)$ espacio topológico. Decimos que $X$ es:
  \begin{enumerate}[(1)]
    \item \wea{Tychonoff} si $\forall u,v \in X$, existen abiertos $\O_u, \O_v \in \t$ tales que
    \[u \in \O_u \setminus \O_v \quad\land\quad
    v \in \O_v \setminus \O_u.\]

    \item \wea{Hausdorff} si $\forall u,v \in X$, existen $\O_u, \O_v \in \t$ tales que
    \[u \in \O_u \quad\land\quad
    v \in \O_v \quad\land\quad
    \O_u \cap \O_v = \varnothing.\]

    \item \wea{Regular} si $X$ es Tychonoff y además para $F \subseteq X$ cerrado y $u \in X \setminus F$ existen $\O, \O_u \in \t$ tales que
    \[F \subseteq \O \quad\land\quad
    u \in \O_u \quad\land\quad
    \O \cap \O_u = \varnothing.\]

    \item \wea{Normal} si $X$ es Tychonoff y si $F_1, F_2 \subseteq X$ son cerrados entonces existen $\O_1,\O_2 \in \t$ tales que
    \[F_1 \subseteq \O_1 \quad\land\quad
    F_2 \subseteq \O_2 \quad\land\quad
    \O_1 \cap \O_2 = \varnothing.\]
  \end{enumerate}
\end{definition}

\begin{proposition}
  $(X,\t)$ es Tychonoff si y solo si $\{v\}$ es cerrado para cualquier $v \in X$.
\end{proposition}
\begin{proof}\
  \begin{itemize}
    \item \boxed{\implies} Supongamos que $(X,\t)$ es Tychonoff.
    Sea $v \in X$, y sea $u \in X \setminus \{v\}$. Luego existe $\O_u \in \t$ tal que $u \in \O_u$ y $v \notin \O_u$. Luego $u \in \O_u \subseteq X \setminus \{v\}$, y como $u$ es arbitrario, $X \setminus \{v\}$ es abierto. Luego $\{v\}$ es cerrado.

    \item \boxed{\impliedby} Supongamos que $\{v\}$ es cerrado para todo $v \in X$. Sean $u,v \in X$ tales que $u \neq v$. Entonces $u \in X \setminus \{v\}$. Como $X \setminus \{v\}$ es abierto, existe $\O_u$ vecindad de $u$ tal que $\O_u \subseteq X \setminus \{v\}$. Por lo tanto $v \notin \O_u$. De la misma forma existe $\O_v$ vecindad tal que $u \notin \O_v$ y $v \in \O_v$.
  \end{itemize}
\end{proof}

\begin{corollary}
  Normal $\subseteq$ Regular $\subseteq$ Hausdorff $\subseteq$ Tychonoff.
\end{corollary}

\begin{definition}
  Sea $(X,d)$ espacio métrico y sean $p \in X$, $C \subseteq X$. La \wea{distancia entre $p$ y $C$} es
  \[\dist(p,C) = \inf\{\dist(p,c) \colon c \in C\}.\]
\end{definition}

\begin{theorem}
  Si $(X,d)$ es un espacio métrico, entonces es normal.
\end{theorem}
\begin{proof}
  Tomemos $F_1,F_2$ cerrados y disjuntos. Definimos
  \begin{align*}
    \O_1 &= \{a \in X \colon \dist(a,F_1) < \dist(a,F_2)\}, \\
    \O_2 &= \{b \in X \colon \dist(b,F_2) < \dist(b,F_1)\}.
  \end{align*}
  Tenemos que $\O_1 \cap \O_2 = \varnothing$. Notemos que si $a \notin F_2$, entonces $\dist(a,F_2)>0$, pues si $\dist(a,F_2) = 0$ entonces existe $(a_n)_{n\in\N} \subseteq F_2$ tal que $\lim_{n\to\oo} a_n = a$, por lo que $a \in \bar{F_2} = F_2$, contradicción. Entonces $F_1 \subseteq \O_1$ y $F_2 \subseteq \O_2$.

  Veamos que $\O_i$ es abierto:
  Sea $a \in \O_1$, por demostrar que $\exists \e>0$ tal que si $d(a,b) < \e$ entonces $b \in \O_1$ (si $b \in B_\e(a)$ entonces $b \in \O$).
  Queremos ver
  \[d(b,F_1) \leq d(b,p_1), \quad \forall p_1 \in F_1.\]
  Notemos que $d(b,F_1) \leq d(a,b) + d(a,p_1)$. Luego $\forall d >0$, $\exists p_1 \in F_1$ tal que
  \[d(a,p_1) \leq \dist(a,F_1) + \d.\]
  Luego $\forall \d>0$,
  \[d(b,F_1) \leq d(a,b) + \dist(a,F_1) + \d.\]
  Tomando $\d \to 0$
  \begin{align*}
    \dist(b,F_1) &\leq d(a,b) + \dist(a,F_1) \\
    &\leq d(a,b) + \dist(a,F_1) - \dist(a,F_2) + \dist(a,F_2).
  \end{align*}
  De la misma manera:
  \[d(a,F_2) \leq d(a,b) + \dist(b,F_2)\]
  Luego
  \begin{align*}
    d(b,F_1) &\leq 2d(a,b) + \dist(a,F_1) - d(a,F_2) + \dist(b,F_2) \\
    \dist(b,F_1) - \dist(b,F_2) &\leq 2d(a,b) + \dist(a,F_1) - \dist(a,F_2).
  \end{align*}
  Tomamos $\e>0$ tal que $2\e < \dist(a,F_2) - \dist(a,F_1)$, y entonces $b \in \O$. Por lo tanto $(X,d)$ es normal.
\end{proof}

\begin{theorem}
  Sea $X$ Tychonoff. $X$ es normal si y solo si $\forall F \subseteq X$ cerrado y $U$ vecindad abierta de $F$, existe $\O$ abierto tal que
  \[F \subseteq \O \subseteq \bar\O \subseteq U.\]
\end{theorem}
\begin{proof}\
  \begin{itemize}
    \item \boxed{\implies} Sea $F$ cerrado y $U$ abierto con $F \subseteq U$. Consideremos $F_2 = X \setminus U$. Tenemos que $F_2$ es cerrado y $F \cap F_2 = \varnothing$. Luego como $X$ es normal, existen $\O_1,\O_2 \in \t$ tales que
    \[F \subseteq \O_1, \quad F_2 \subseteq \O_2, \quad \O_1 \cap \O_2 = \varnothing.\]
    Veamos que $\O_1$ cumple. Ta tenemos que $F \subseteq \O_1$, además $\O_1 \cap \O_2 = \varnothing$ y $\O_2 \supseteq X \setminus U$. Por lo tanto $X \setminus \O_2 \subseteq U$. Además $\O_1 \subseteq X \setminus \O_2$
    \[\implies \bar{\O_1} \subseteq \bar{X \setminus \O_2} = X \setminus \O_2 \subseteq U.\]
    Por lo tanto $F \subseteq \O_1 \subseteq \bar{\O_1} \subseteq U$.

    \item \boxed{\impliedby} Sean $F_1, F_2$ cerrados tales que $F_1 \cap F_2 = \varnothing$. Sea $U = X \setminus F_2$, entonces $F_1 \subseteq U$.
    Por propiedad, existe $\O \in \t$ tal que
    \[F_1 \subseteq \O \subseteq \bar\O \subseteq U = X \setminus F_2.\]
    Tomamos $\O_2 = X \setminus \bar\O \in \t$. Entonces $F_2 \subseteq \O_2$. Similar para $F_1$.
  \end{itemize}
\end{proof}

\begin{definition}
  Un espacio topológico $(X,\t)$ se dice \wea{primero contable} si $\forall p \in X$ existe una base de abiertos en $p$ tal que la base es numerable.
\end{definition}

\begin{definition}
  $(X,\t)$ se dice \wea{segundo contable} si existe una base de $\t$ numerable.
\end{definition}

\begin{example}
  Todo espacio métrico $(X,d)$ es primero contable. Todo punto $p \in X$ tiene base $\{B_r(p) \colon p \in \Q_{>0}\}$.
\end{example}

\begin{definition}
  Sea $(X,\t)$ espacio topológico. Decimos que $A \subseteq X$ es \wea{denso} si $\bar{A} = X$.
\end{definition}

\begin{definition}
  Un espacio topológico $(X,\t)$ se dice \wea{separable} si existe un subconjunto denso numerable.
\end{definition}

\begin{theorem}
  Sea $(X,d)$ espacio métrico. $X$ es segundo contable si y solo si $X$ es separable.
\end{theorem}
\begin{proof}\
  \begin{itemize}
    \item \boxed{\impliedby} Tomamos $E = \{e_i\}_{i\in\N}$ denso. Tomamos $$\B = \{B_r(e_i) \colon r \in \Q_{>0}, e_i \in E\}.$$
    Veamos que es base:
    \begin{enumerate}[(i)]
      \item $\bigcup_{B\in\B} B = X$: trivial %% demostrar
      \item $B_1, B_2\in \B \implies \exists B_3 \subseteq B_1 \cap B_2$:
      
      Sean $B_1, B_2$. Tenemos que< $B_1 \cap B_2 \neq \varnothing$. Sea $p \in B_1 \cap B_2$, luego existen $r_1,r_2$ tales que $B_{r_i}(p) \subseteq B_i$. Como $E$ es denso, existe $e_i \in B_{r_i}(p)$, y como $B_{r_i}(p)$ son abiertos, existe $\e$ tal que $B_\e(e_i) \subseteq B_{r_i}(p)$. Basta con tomar $\e$ racionales.
    \end{enumerate}

    \item \boxed{\implies} Sea $\{B_i\}_{i\in\N}$ base de abiertos numerable. Basta tomar $e_i \in B_i$, y considerar $E = \{e_i\}_{i\in\N}$.
  \end{itemize}
\end{proof}

\begin{theorem}
  Si $(X,\t)$ es un espacio topológico primero contable, entonces $a \in \bar{E}$ si y solo si $\exists (a_n)_{n\in\N} \subseteq E$ tal que $\lim_{n\to\oo} a_n = a$.
\end{theorem}
\begin{proof}\
  \begin{itemize}
    \item \boxed{\implies} Sea $p \in \bar{E}$. Tomamos $\{B_n\}$ base de abiertos de $p$ numerable. Notemos que, como $p \in \bar{E}$, entonces $B_1 \cap E \neq \varnothing$. Tomamos $p_1 \in B_1 \cap E$. Luego
    \[p_i \in E \cap \bigcap_{j\leq i} B_i \neq \varnothing.\]
    Este conjunto es no vacío porque es abierto y $p \in B_1 \cap \dots \cap B_i$.
    Ahora veamos que $p_n \to p$.

    Tomamos $B$ abierto con $p \in B$. Como $\{B_i\}$ es base, existe $i_0$ tal que $B_{i_0} \subseteq B$. Luego para todo $i \geq i_0$
    \[p_i \in B_i \cap \dots \cap B_{i_0} \cap \dots \cap B_1 \cap E.\]
    Luego $p_i \in B_{i_0}$. Para todo $i \geq i_0$, $p_i \in B$.
  \end{itemize}
\end{proof}

\begin{theorem}
  [Urysohn]
  Si $(X,\t)$ es normal y segundo contable, entonces $(X,\t)$ es metrizable.
\end{theorem}

\newpage
\fecha{Clase 3. Lunes 16 de Marzo}
\subsection{Funciones Continuas}

\begin{definition}
  Si $(X,d)$ e $(Y,\rho)$ son espacios métricos, un función $f \colon X \to Y$ se dice \wea{continua en $a \in X$} si $\forall \e>0, \exists \d>0$ tal que $$d(a,b) < \d \implies \rho(f(a),f(b)) < \e.$$
  La función se dice \wea{continua} si es continua en todo $a \in X$.
\end{definition}

\begin{obs}
  Esto es equivalente a que $\forall \e>0, \exists \d>0$ tal que $$b \in B_\d(a) \implies f(b) \in B_\e(f(a)).$$
\end{obs}

\begin{obs}
  Esto es equivalente a que si $\O$ es una vecindad de $f(a)$, entonces existe $U$ vecindad de $a$ tal que $f(U) \subseteq \O$.
\end{obs}

\begin{obs}
  Hay dos definiciones equivalentes:
  \begin{enumerate}[(1)]
    \item Si $\O$ es abierto en $Y$, entonces $f^{-1}(\O)$ es abierto en $X$.
    \item $f$ es continua en $a$ si $\forall a_n \to a \implies f(a_n) \to f(a)$.
  \end{enumerate}
\end{obs}

\begin{definition}
  Sean $(X,\t_X)$ e $(Y,\t_Y)$ espacios topológicos. Decimos que $f \colon X \to Y$ es \wea{continua en $a \in X$} si $\forall \O \in \t_Y$ vecindad de $f(a)$, existe $U \in \t_X$ vecindad de $a$ tal que $f(U) \subseteq \O$.
\end{definition}

\begin{proposition}
  Sean $(X,\t_X), (Y,\t_Y)$ espacios topológicos. Una función $f \colon X \to Y$ es continua si y solo si $\forall \O \in \t_Y$, $f^{-1}(\O) \in \t_X$.
\end{proposition}
\begin{proof}
  \
  \begin{itemize}
    \item [\boxed{\implies}]
    Sea $\O \in \t_Y$ no vacío. Tenemos que $f^{-1}(\O) \neq \varnothing$, por lo que existe $a \in X$ tal que $f(a) \in \O$.
    Por la definición, existe $U_a \in \t_X$ tal que $a \in U_a$ y $f(U_a) \subseteq \O$. $$\implies U_a \subseteq f^{-1}(\O).$$
    Tenemos entonces que $\forall a \in f^{-1}(\O), \exists U_a \in \t_X$ tal que $a \in U_a \subseteq f^{-1}(\O)$. Por lo tanto $f^{-1}(\O) \in \t_X$.

    \item [\boxed{\impliedby}]
    Sea $a \in X$ y $\O_a$ vecindad de $f(a)$. Entonces $f^{-1}(\O_a) \in \t_X$. Además por construcción, $a \in f^{-1}(\O_a)$ y entonces $\exists U_a \in \t_X$ tal que $a \in U_a \subseteq f^{-1}(\O_a)$. Por lo tanto $f(U_a) \subseteq \O_a$.
  \end{itemize}
\end{proof}

¿Cuándo podemos caracterizar la continuidad a través de sucesiones?
Tomemos $(a_n)_{n\in\N}$ tal que $a_n \to a$. Esto es si y solo si
\[\tag{$*$}
\forall U_a \in \t_X [a \in U_a \implies \exists n_0 \in \N (n \geq n_0 \implies a_n \in U_a)]
\]

%% me faltó algo acá, lo completo después

Tomemos $\O$ vecindad de $f(a)$. Tomemos $U = f^{-1}(\O)$. Como $f$ es continua, $U$ es abierto, y por $(*)$ existe $n_0$ tal que $n \geq n_0 \implies a_n \in U = f^{-1}(\O)$, y entonces $f(a_n) \in \O$
$$\implies f(a_n) \to f(a).$$

¿Se cumplirá que si $$a_n \to a \implies f(a_n) \to f(a)$$ entonces $f$ es continua?
Habría que demostrar
\[\tag{$**$}
\forall \O \in \t_Y, \quad f^{-1}(\O) \in \t_x.
\]
Si fueran espacios métricos, $f^{-1}(B_{1/n}(f(a)))$ no es abierto para infinitos valores de $n$.
Entonces
\[\exists a_n \in f^{-1}(B_{1/n}(f(a))) \colon B_{r_n}(a) \not\subseteq f^{-1}(B_{1/n}(f(a)))\]
\[\implies \exists b_n \in B_{r^n}(a) \ \land \ f(b_n) \notin B_{1/n}(f(a)).\]

\begin{proposition}
  \
  \begin{enumerate}[(1)]
    \item Si $f \colon X \to Y$ es continua, entonces \[\forall (a_n)_{n\in\N} \cor{\lim_{n\to\oo} a_n = a \implies \lim_{n\to\oo} f(a_n) = f(a)}\]
    \item Si $X$ e $Y$ son primero contables, entonces
    \[\forall (a_n)_{n\in\N}\cor{\lim_{n\to\oo} a_n = a \implies \lim_{n\to\oo} f(a_n) = a} \implies f \text{ es continua}.\]
  \end{enumerate}
\end{proposition}
\begin{proof}
  \
  \begin{enumerate}[(1)]
    \item \checkmark %% meter demostración acá
    \item Tomemos $\B = \{\O_n\}_{n\in\N}$ base de vecindades de $f(a)$ y $\{U_n\}_{n\in\N}$ base de abiertos de $a$. Definimos $$\tilde\O_n = \O_1 \cap \dots \cap \O_n.$$ Son abiertos (intersección finita) y $\tilde\O_{n+1} \subseteq \tilde\O_n$. Definimos $$\tilde U_n = U_1 \cap \dots \cap U_n.$$ Queremos demostrar
    \[\forall \O \in \t_Y, \quad f^{-1}(\O) \in \t_X.\]
    Por contradicción, supongamos
    \[\exists \O_0 \in \t_Y \colon f^{-1}(\O_0) \notin \t_X.\]
    Luego existe $a_0 \in f^{-1}(\O_0)$ que no pertenece al interior de $f^{-1}(\O)$.
    \[\implies \forall U \in \t_X \cor{a_0 \in U \implies U \not\subseteq \bar{f^{-1}(\O_0)}}\]
    Tomamos base contable encajada de $a_0$, luego $(a_n)_{n\in\N}$ y $a_n \in U_n \cap (f^{-1}(\O))^c$
    \[\implies \lim_{n\to\oo} a_n = a \implies \lim_{n\to\oo} f(a_n) = f(a).\]
    Pero $f(a_n) \notin \O$ y $f(a) \in \O$. Por lo tanto $f(a_n) \not\to f(a)$.
  \end{enumerate}
\end{proof}

\begin{proposition}
  Si $(X,\t_X), (Y,\t_Y), (Z,\t_Z)$ son espacios topológicos, y las funciones $f \colon X \to Y$ y $g \olon Y \o Z$ son continuas, entonces $g \circ f$ es continua.
\end{proposition}

Pregunta: Dados $(X,\t), \{(X_\a,\t_\a)\}_{\a\in I}$ espacios topológicos y $f_\a \colon X \to X_\a$ funciones continuas $\forall \a \in I$, ¿Cuál es la topología con ``menos'' abiertos (más débil) que preserva la continuidad de estas funciones?

\begin{definition}
  Si $\t_1$ y $\t_2$ son topologías en $X$, decimos que $\t_2$ es \wea{más débil} que $\t_1$ si $$\t_2 \subseteq \t_1.$$
\end{definition}

Si tomamos $$\t_\a' = \{f_\a^{-1}(\O_\a) \colon \O_a \in \t_\a\}$$ entonces $\t_\a'$ tiene que estar incluido en esa topología débil.

\begin{proposition}[o definición no se]
  Cualquier familia de conjuntos genera una base de abiertos y la familia se llama \wea{subbase de la topología}.
\end{proposition}
\begin{proof}
  $$\Fcal = \{C_\a \colon \a \in A\}$$
  Si $$\B = \set{\bigcap_{i=1}^n C_i \colon C_i \in \Fcal} \cup \{X\}$$
  es fácil ver que $\B$ cumple las condiciones de la proposición sobre bases, pues
  \[\br{\bigcap_{i=1}^n C_i \cap \bigcap_{j=1}^m \tilde C_j} \in \B \quad\land\quad C_i \in \B.\]
\end{proof}
%% reescribir todo esto, buscar definición de subbase en el libro (no lo he leido)

\begin{definition}
  Sean $(X,\t_X), (Y,\t_Y)$ espacios topológicos. Una función $f \colon X \to Y$ se dice \wea{homeomorfismo} si es continua, biyectiva y tiene unversa continua.
\end{definition}

\begin{obs}
  Los homeomorfismos definen una clase de equivalencia entre espacios topológicos.
\end{obs}

\newpage
\subsection{Compacidad}

\begin{definition}
  Sea $X$ un conjunto. Un \wea{cubrimiento} es una familia $\{\O_\a\}_\a$ tal que $\ds X = \bigcup_\a \O_\a$. Si $(X,\t)$ es un espacio topológico, decimos que $\{\O_\a\}_\a$ es un \wea{cubrimiento abierto} ( o \wea{cubrimiento por abiertos}) si los $\O_\a$ son abiertos.
\end{definition}

\begin{definition}
  Un espacio métrico $(X,d)$ se dice \wea{totalmente acotado} si $\forall \e>0, \exists \{x_i\}_{i=1}^n$ tal que $$X = \bigcup_{i=1}^n B_\e(x_i).$$
\end{definition}

\begin{definition}
  Un espacio métrico $(X,d)$ se dice \wea{compacto} si Para cualquier cubrimiento abierto de $X$ existe un subcubrimiento finito.
\end{definition}

\begin{theorem}
  En espacios métricos son equivalentes:
  \begin{enumerate}[(1)]
    \item Todo cubrimiento abierto tiene subcubrimiento finito;
    \item Toda sucesión tiene subsucesión convergente.
    \item Completo y totalmente acotado;
  \end{enumerate}
\end{theorem}

Todas estas propiedades sirven como definición de espacio métrico compacto.

\fecha{Clase 4. Miércoles 18 de Marzo}

\begin{obs}
  Si $(X,\t)$ es un espacio topológico, no tenemos una forma de interpretar la definición (iii): No sabemos que es una sucesión de Cauchy sin métrica, y por lo tanto no tenemos completitud. Tampoco sabemos que es ser totalmente acotado en espacios topológicos generales. Por estas razones consideramos (i) como definición de compacidad para el caso general.
\end{obs}

\begin{proposition}
  Sea $(X,\t)$ espacio topológico. Si $\Fcal = \{F_\a\}_{\a \in A}$ son cerrados tales que $$\forall \Ucal \subseteq \Fcal \cor{\#\Ucal < \oo \implies \bigcap \Ucal \neq \varnothing}$$ entonces $$\bigcap \Fcal \neq \varnothing.$$
\end{proposition}

Esto escrito con abiertos es

\[\bigcup_{\a \in A} F_\a^c = X \iff \br{\bigcap_{\a \in A} F_\a = \varnothing \implies \exists F_1,\dots,F_n \cor{F_1 \cap \dots \cap F_n = \varnothing}}\]

\begin{proposition}
  Si $(X,\t)$ es compacto y $F \subseteq X$ es cerrado, entonces $F$ es compacto con la topología inducida.
\end{proposition}
\begin{proof}
  Recordar que $\O_F$ es bierto con la topologí inducida, $\O_F = \O \cap F$ con $\O$ abierto en $X$. Si tomamos un cubrimiento abierto de $F$:
  \[F = \bigcup_{\a \in A} (\O_\a \cap F)\]
  con $\O_\a$ abiertos en $X$. Como $F$ es cerrado, entonces $X \setminus F$ es abierto y $\{\O_\a\}_{\a \in A} \cup \{X \setminus F\}$ cubre $X$. Como $X$ es compacto, existen $\O_1, \dots, \O_n$ tales que $\ds \bigcup_{i=1}^n \O_i = X$. Luego
  \[F = \bigcup_{i=1}^n \O_i \cap F = \bigcup_{i=1}^n (\O_i \cap F).\]
\end{proof}

\begin{theorem}
  Si $(X,\t)$ es compacto y Hausdorff, y $F \subseteq X$ es compacto, entonces $F$ es cerrado.
\end{theorem}
\begin{proof}
  Notemos que $F$ es cerrado si y solo si $X \setminus F$ es abierto. Si $X \setminus F = \varnothing$, entonces $F=X$ y  es cerrado. Si $X \setminus F \neq \varnothing$, tomamos $a \in X \setminus F$ y $f_\a \in F$. Como $X$ es Hausdorff, entonces existen abiertos disjuntos $\O_a, U_\a$ tales que $a \in \O_\a$ y $f_\a \in U_\a$.
  Sea $\{U_\a \colon f_\a \in F\}$ cubrimiento abierto de $F$. Como $F$ es compacto, tenemos $U_1,\dots,U_n$ que cubren $F$.
  Tomamos $\O = \O_1 \cap \dots \cap \O_n$ tales que $a \in \O_i$ y $\O_i \cap U_i = \varnothing$. Tenemos que $\O$ es abierto. Luego
  \[\O \cap F = \O \cap \bigcup_{i=1}^n U_i = \bigcup_{i=1}^n (\O \cap U_i) = \varnothing\]
  \[\implies \O \subseteq X \setminus F.\]
  Luego $X \setminus F$ es abierto y entonces $F$ es cerrado.
\end{proof}

\begin{definition}
  $(X,\t)$ se dice \wea{secuencialmente compacto} si para toda sucesión existe una subsucesión convergente.
\end{definition}

\begin{proposition}
  Sea $(X,\t)$ primero contable y compacto. Entonces $X$ es secuencialmente compacto.
\end{proposition}

\begin{definition}
  $(X,\t)$ se dice \wea{Lindelöf} si toso cubrimiento abierto tiene un cubcubrimiento contable.
\end{definition}

\begin{obs}
  Compacto $\implies$ Lindelöf.
\end{obs}

\begin{lemma}
  Segundo contable $\implies$ Lindelöf.
\end{lemma}
\begin{proof}
  Tomemos $\{\O_i\}_{i\in\N}$ base de abiertos numerable y $\{U_\a\}_{\a \in A}$ cubrimiento. Tomemos $U_i$ tal que $\O_i \subseteq U_i$. Luego
  \[\bigcup_{i\in\N} U_i \supseteq \bigcup_{i\in\N} \O_i = X.\]
\end{proof}

\begin{theorem}
  Sea $(X,\t)$ segundo contable. $X$ es compacto si y solo si $X$ es secuencialmente compacto.
\end{theorem}

\fecha{Clase 5. Lunes 13 de Marzo}

\begin{theorem}
  Si $(X,\t)$ es compact y Hausdorff, entonces $(X,\t)$ es normal.
\end{theorem}
\begin{proof}
  Royden, p. 235, Teorema 18.
\end{proof}

\begin{theorem}
  Sea $(X,\t)$ compacto y $f \colon X \to Y$ continua. Entonces $f(X)$ es compacto.
\end{theorem}
\begin{proof}
  Royden, p. 235, Proposición 20.
\end{proof}

\begin{corollary}
  Si $f \colon X \to \R$ es continua y $X$ es compacto, entonces $f$ alcanza un máximo y mínimo.
\end{corollary}

\begin{corollary}
  Sea $X$ compacto e $Y$ Hausdorff. Si $f \colon X \to Y$ es una función continua y biyectiva, entonces $f$ es homeomorfismo.
\end{corollary}

\newpage
\subsection{Conexidad}

\begin{definition}
  Un espacio topológico $(X,\t)$ se dice \wea{conexo} si
  \[\forall \O_1, \O_2 \in \t \cor{
    (\O_1 \cap \O_2 \ \land\ \O_1 \cup \O_2 = X) \implies \O_1, \O_2 \in \{\varnothing, X\}
  }.\]
  Si $E \subseteq X$, decimos que $E$ es \wea{conexo} si es conexo con la topología inducida.
\end{definition}

\begin{theorem}
  Si $X$ es conexo y $f \colon X \to Y$ es continua, entonces $f(X)$ es conexo.
\end{theorem}

\begin{exercise}
  En $\R$ los conexos son intervalos.
\end{exercise}

\begin{corollary}
  Si $X$ es conexo y $f \colon X \to \R$ continua, entonces $f(X)$ es un intervalo.
\end{corollary}

\begin{definition}
  $X$ es un espacio topológico tal que $\forall f \colon X \to \R$, $f(X)$ es intervalo, se dice que $X$ cumple la \wea{propiedad del valor intermedio} (abreviado como \wea{PVI}).
\end{definition}

\begin{theorem}
  $X$ cumple PVI si y solo si $X$ es conexo.
\end{theorem}
\begin{proof}
  Supongamos que $X$ no es conexo. Veamos que no cumple PVI.
  Si no es conexo, entonces existen $\O_1,\O_2$ abiertos tales que
  \[\O_1 \neq \varnothing \neq \O_2 \quad\land\quad
  \O_1 \cap \O_2 = \varnothing \quad\land\quad
  X = \O_1 \cup \O_2.\]
  Consideremos la función
  \[f(x) = \begin{cases}
    1, &x \in \O_1 \\
    0, &x \in \O_2
  \end{cases}\]
  Entonces
  \[f^{-1}((a,b)) = \begin{cases}
    \O_1, &\text{si } 1 \in (a,b), 0 \notin (a,b), \\
    \O_2, &\text{si } 0 \in (a,b), 1 \notin (a,b), \\
    X, &\text{si } 0,1 \in (a,b), \\
    \varnothing, &\text{si } 0,1 \notin (a,b).
  \end{cases}\]
  $f$ continua, entonces $f(X)$ no es intervalo.
\end{proof}

\begin{lemma}
  [\wea{Urysohn}]
  Sea $X$ normal y sean $A,B \subseteq X$ cerrados disjuntos. Sea $[a,b]$ cerrado, entonces existe $f \colon X \to [a,b]$ continua tal que
  \[f(A) = \{a\} \quad\land\quad f(B) = \{b\}.\]
\end{lemma}

\begin{obs}
  Si $(X,d)$ es un espacio métrico, entonces tenemos
  \[f(x) = \frac{b \cdot \dist(x,A)}{\dist(x,A) + \dist(x,B)} + \frac{a \cdot \dist(x,B)}{\dist(x,A) + \dist(x,B)}.\]
\end{obs}

\begin{obs}
  Si un espacio que cumple que $\forall A,B$ cerrados disjuntos existe $f \colon X \to [a,b]$ continua tal que $f(A) = \{a\}$ y $f(B) = \{b\}$, entonces es normal, pues si $A,B$ son cerrados disjuntos, tomamos
  \[\O_1 = f^{-1}([a,(b-1)/3)) \supseteq A \quad\land\quad \O_2 = f^{-1}((2(b-a)/3, b]) \supseteq B.\]
\end{obs}

\begin{definition}
  Sea $X$ un espacio topológico y $\Lambda \subseteq \R$. Una familia $\{\O_\lambda\}_{\lambda\in\Lambda}$ se dice \wea{noramlmente ascendente} si \[\forall \lambda_1,\lambda_2 \in \Lambda \cor{\lambda_1 < \lambda_2 \implies \bar{\O_{\lambda_1}} \subseteq \O_{\lambda_2}}.\]
\end{definition}

\begin{lemma}\label{lemma Urysohn proof 1}
  Sea $X$ un espacio topológico y $\Lambda$ un conjunto denso en $(a,b)$. Sea $\{\O_\lambda\}_{\lambda\in\Lambda}$ familia normalmente ascendente de abiertos. Si tomamos
  \[f(x) = \begin{cases}
    b, &si \ X \setminus \bigcup_{\lambda\in\Lambda} \O_\lambda \\
    \inf\{\lambda\in\Lambda \colon x \in \O_\lambda\}, &e.o.c.
  \end{cases}\]
  entonces $f$ es continua.
\end{lemma}
\begin{proof}
  Tomemos $y_0 \in X$ y $f(y_0) = r_0$. Por demostrar que $\forall \e>0$, $f^{-1}(r_0-\e, r_0+\e)$ es abierto.
  Esto es si y solo si $\forall \e>0$, $\exists U$ abierto en $X$ tal que
  \[y_0 \in U \subseteq f^{-1}(r_0-\e,r_0+\e).\]
  Asumamos primero que $r_0 \in (a,b)$. Existen $\lam_1,\lam_2 \in \Lam$ tales que
  \[a < \lam_1 < r_0 < \lam_2 < b\]
  Luego existen $\O_{\lam_1}, \O_{\lam_2} \in \{\O_\lam\}$ tales que
  \[\bar{\O_{\lam_1}} \subseteq \O_{\lam_2}.\]
  Tomemos
  \[U = \O_{\lam_2} \setminus \bar{\O_{\lam_1}}.\]
  Tenemos que $y_0 \in U$ porque $f(y_0) = r_0$ y entonces $\lam_1 < f(y_0) < \lam_2$.
\end{proof}

\begin{lemma}\label{lemma Urysohn proof 2}
  Sea $X$ normal. Si $F$ es cerrado y $U$ es vecindad abierta de $F$, entonces existe una familia $\{\O_\lambda\}_{\lambda\in\Lambda}$ normalmente ascendente de abiertos tal que
  \[F \subseteq \O_\lambda \subseteq \bar{\O_\lambda} \subseteq U.\]
\end{lemma}

\begin{proof}[\dem{Urysohn}]%% Urysohn, mover esto al enunciado, y poner los lemas antes.
  Tenemos $A,B$ cerrados. En el Lema \ref{lemma Urysohn proof 2} tomamos $F=A$ y $U = X \setminus B$.
  Si $A \cap B = \varnothing$, entonces $A \subseteq U$, luego el lema \ref{lemma Urysohn proof 2} entrega $\{\O_\lambda\}_{\lambda\in\Lambda}$ y construimos $f$ como en el lema \ref{lemma Urysohn proof 1}.
\end{proof}

\fecha{Clase 6. Miércoles 25 de Marzo}

\begin{lemma}
  Sea $\{f_n\}_{n\in\N}$ con $f_n \colon X \to [a,b]$ continuas, tal que
  \[f(x) = \sum_{n\in\N} f_n(x)\]
  es uniformemente convergente. Entonces $f$ es continua.
\end{lemma}
\begin{proof}
  Dado $y_0 \in X$, $\forall \e>0$, $\exists U \subseteq X$ vecindad de $y_0$ tal que
  \[y \in U \implies \abs{f(y) - f(y_0)} < \e.\]
  Como $\ds \sum_{n=1}^\oo f_n(y)$ es uniformemente convergente, existe $n_0$ tal que
  \[\abs{\sum_{n=n_0}^\oo f_n(x)} < \frac{\e}{3}\]
  \[\implies \abs{\sum_{n=1}^{n_0-1} f_n(y) - f_n(y_0)} \leq \frac{\e}.\]
  Como cada $f_n$ es continua, $\exists U_n$ tal que
  \[y\in U_n \implies \abs{f_n(y) - f_n(y_0)} < \frac{\e}.\]
  Tomando $U = U_1 \cap \dots \cap U_{n_0-1}$, tenemos que
  \[y \in U \implies \abs{\sum_{n=1}^{n_0-1} f_n(y) - f_n(y_0)} \leq \frac{\e}{3}.\]
\end{proof}

\begin{theorem}
  [Tietze]
  Si $X$ es normal, $F \subseteq X$ es cerrado y $f \colon F \to [a,b]$ continua, entonces existe $g \colon X \to [a,b]$ tal que $g\vert_F = f$.
\end{theorem}
\begin{proof}
  Suponemos que $[a,b] = [-1,1]$. La idea es construir funciones $g_n \colon X \to [-1,1]$ continuas tales que $\abs{g_n(x)} \leq (2/3)^n$ y para $x \in F$
  \[\abs{f(x) - \sum_{i=1}^n g_i(x)} \leq \br{\frac{2}{3}}^n.\]
  Entonces la extensión de $f$ es la función
  \[g(x) = \sum_{i=1}^\oo g_i(x).\]
  \begin{claim}
    $\forall K>0$ y $h \colon F \to \R$ tal que $\abs{h}\leq K$, $\exists g$ continua tal que $\abs g < 2K/3$ en $X$ y $\abs{h-g} < 2K/3$.
  \end{claim}
  \begin{obs}
    Si la afirmación se cumple:
    \begin{itemize}
      \item Para $n=1$, tomamos $h=f$ y $K=1$, obtenemos $g_1$.
      \item Para $n=2$, tomamos $h = f-g_1$ y $K = 2/3$, obtenemos $g_2$ con $\abs{g_2} < 1/3$.
      \item Para $n$ cualquiera, $h = f - g_1 - \dots - g_{n-1}$ y $K = (2/3)^{n-1}$. Luego existe $g_n$ tal que
      \[\abs{g_n} \leq \br{\frac{2}{3}}^n \quad\land\quad \abs{f - \sum_{i=1}^n g_i} < \br{\frac{2}{3}}^n.\]
    \end{itemize}
  \end{obs}
  Para demostrar la afirmación, tomamos
  \[A = \{x \in F \colon h(x) \leq -\frac{K}{3}\},\]
  \[B = \{x \in F \colon h(x) \geq \frac{K}{3}\}.\]
  Por el Lema de Urysohn, $\exists g \colon X \to [-K/3,K/3]$ continua tal que
  \[g(x) = \begin{cases}
    -K/3, &x \in A \\
    K/3, &x \in B.
  \end{cases}\]
  En $A$ tenemos que
  \[-K \leq h(x) \leq -\frac{K}{3},\]
  \[-\frac{2K}{3} \leq h(x) - g(x) \leq 0.\]
  En $B$ tenemos que
  \[\frac K3 \leq h(x) \leq K,\]
  \[0 \leq h(x) - g(x) \leq \frac{2K}{3}.\]
  Entonces este es el $g$ de la afirmación.
\end{proof}

\begin{theorem}
  [Teorema de Metrización de Urysohn]
  Sea $X$ segundo contable. Entonces $X$ es normal si y solo si $X$ es metrizable.
\end{theorem}
\begin{proof}
  Ya sabemos que metrizable implica normal, así que supongamos que $X$ es segundo contable y normal, y mostremos que es metrizable.

  Sea $\{U_n\}_{n\in\N}$ base de abiertos numerables. Sea
  \[A = \{(n,m) \in \N \times \N \colon \bar{U_n} \subseteq U_m\}.\]
  Por el Lema de Urysohn, $\forall (n,m) \in A, \exists f_{n,m} \colon X \to [0,1]$ tal que $\abs{f_{n,m}} < 1$ y
  \[f_{n,m}(x) = \begin{cases}
    0, &x \in \bar{U_n} \\
    1, &x \in X \setminus U_m.
  \end{cases}\]
  Tomemos
  \[d(x,y) = \sum_{(n,m) \in A} \frac{\abs{f_{n,m}(x) - f_{n,m}(y)}}{2^{n+m}}.\]
  \begin{claim}
    $d$ es métrica en $X$.
  \end{claim}
  Veamos que
  \begin{itemize}
    \item \boxed{d(x,y) = 0 \iff x=y} Notemos que si $x=y$ entonces
    \[d(x,x) = \sum_{(n,m) \in A} \frac{\abs{f_{n,m}(x) - f_{n,m}(x)}}{2^{n+m}} = \sum_{(n,m) \in A} \frac{0}{2^{n+m}} = 0.\]
    Si $x \neq y$, entonces $\{x\} \subseteq X \setminus \{y\} \eqcolon U$. Como $X$ es normal, entonces existe $\O$ tal que
    \[
    \{x\} \subseteq \O \subseteq \bar\O \subseteq U.\]
    Tomamos $U_n$ vecindad de $x$ en $\O$ y $U_m = U$ (???)
    \[(n,m) \in A \implies f_{n,m}(x) = 0 \ \land \ f_{n,m}(y) = 1.\]
    \[\implies d(x,y) > 0.\]
    \item \boxed{d(x,y) = d(y,x)} Trivial.
    \item \boxed{d(x,y) \leq d(x,z) + d(z,y)} Trivial
  \end{itemize}
  \begin{claim}
    La topología de $(X,d)$ coincide con la topología original. Es decir, $\id \colon (X,\t) \to (X,d)$ es homeomorfismo.
  \end{claim}
  Notemos que $d(\cdot, y)$ es continua en $(X,\t)$, y entonces $\id$ es continua. Veamos que es cerrada. Sea $F$ cerrado en $(X,\t)$, queremos ver que $\id(F) = F$ es cerrado en $(X,d)$.

  Tomemos $x \notin F$, entonces existe $\O$ abierto tal que
  \[F \subseteq \O \subseteq \bar\O \subseteq X \setminus \{x\}.\]
  Por demostrar que
  \[\exists \e>0 \cor{y \in F \implies d(x,y) > \e}\]
\end{proof}

\newpage
\fecha{Clase 7. Lunes 30 de Marzo}
\subsection{Teoremas Importantes de Espacios Métricos}
\subsubsection{Teorema de Stone--Weierstrass}

\begin{definition}
  Sea $(X,\tau)$ un espacio topológico.
  $C(X) = C(X;R) = \{f \colon X \to \R \mid f \text{ continua}\}$.
\end{definition}

\begin{definition}
  $$d(f,g) = \sup_{u\in X} \abs{f(u) - g(u)}.$$
  Si $X$ es compacto, entonces $$d(f,g) = \max_{u \in X} \abs{f(u) - g(u)}.$$
  La topología inducida por esta métrica se llama \wea{topología uniforme}.
\end{definition}

\begin{obs}
  $C(X)$ es un $\R$--espacio vectorial.
\end{obs}

\begin{definition}
  Decimos que $A \subseteq C(X)$ es una \wea{álgebra de funciones} si es un subespacio vectorial y si $\forall f,g \in A$, $fg \in A$.
\end{definition}

\begin{definition}
  Se dice que $A \subseteq C(X)$ \wea{separa puntos} si $$\forall u,v \in X, \exists f \in A \colon f(u) \neq f(v).$$
\end{definition}

\begin{theorem}[Stone--Weierstrass]
  Si $X$ es compacto y Hausdorff, y $A$ es un álgebra que separa puntos y que contiene las constantes, entonces $A$ es denso en $C(X)$ con la topología uniforme.
\end{theorem}

\begin{obs}
  Si $X$ es normal, el lema de Urysohn garantiza que $C(X)$ separa puntos. En general, no sabemos si $C(X)$ separa puntos.
\end{obs}

\begin{lemma}
  Sea $X$ compacto Hausdorff y $A$ un álgebra de funciones que separa puntos y contiene las constantes.
  Entonces para todo $F$ cerrado y todo $x_0 \in X \setminus F$, existe $U$ vecindad de $x_0$ tal que
  \[U \cap F = \varnothing \quad\land\quad
  \forall \e > 0, \exists h \in A \cor{
    h\vert_U < \e \quad\land\quad
    h\vert_F > 1-\e \quad\land\quad
    0 \leq h \leq 1
  }\]
\end{lemma}

\begin{lemma}
  Sean $X$ compacto Hausdorff y $A$ un álgebra de funciones que separa puntos y contiene las constantes. Entonces
  \[\forall F_1,F_2 \in \tau^c, \forall \e>0 \cor{
    F_1 \cap F_2 = \varnothing \implies \exists h \in A \br{
      h\vert_{F_1} < \e \quad\land\quad h\vert_{F_2} > 1-\e \quad\land\quad 0 \leq h \leq 1.
    }
  }\]
\end{lemma}

\begin{proof}
  [\dem{Stone--Weierstrass}]
  \begin{claim}
    Basta considerar $f \in C(X)$ tales que $0 \leq f(x) \leq 1$.
  \end{claim}
  En efecto, para $f$ cualquiera, si $c = \max_{u \in X} \abs{f(u)}$ ($X$ es compacto), tenemos que $c \in A$.
  Si $f \in A$, entonces $f+c \geq 0$, $f+c \in A$, y entonces
  \[\frac{f+c}{\abs{f+c}} \in A, \quad 0 \leq \frac{f+c}{\abs{f+c}} \leq 1.\]
  Tomemos $n \in \N$, $f \in C(X)$ con $0 \leq f \leq 1$. Definimos
  \begin{align*}
    C_j &= \set{u \in X \colon f(u) \leq \frac{j-1}{n}}, \\
    B_j &= \set{u \in X \colon f(u) \geq \frac{j}{n}}.
  \end{align*}
  Por el lema anterior, existen $g_j \in A$ tal que $g_j < 1/n$ en $C_j$, $g_j > 1 - 1/n$ en $B_j$ y $0 \leq g_j \leq 1$.
  Tomando
  \[g(x) = \frac{1}{n} \sum_{j=1}^n g_j \in A\]
  Tendríamos
  \[\tag{$*$} \max_{x \in X} \abs{f(x) - g(x)} < \frac{3}{n}.\]
  Demostremos $(*)$. Tomamos $k \in \{1,\dots,n\}$, y supongamos que $0 \leq f(x) \leq k/n$.
  Me rindo con la demostración.
\end{proof}

\fecha{Clase 8. Miércoles 1 de Abril}

\begin{theorem}
  [Riesz]
  Sea $X$ compacto y Hausdorff. Entonces $C(X)$ es separable si y solo si $X$ es metrizable.
\end{theorem}
\begin{proof}\
  \begin{itemize}
    \item \boxed{\implies} Supongamos que $X$ es separable. Sabemos que compacto y Hausdorff implica normal, entonces si $X$ es segundo contable, obtenemos por Urysohn que es metrizable.

    Existe $\{g_n\}_{n\in\N}$ denso en $C(X)$. Tomemos
    \[U_n = \{x \in X \colon g_n > 1/2\}.\]
    Veamos que $\{U_n\}_{n\in\N}$ es base de la topología de $X$.
    Sea $\O \subseteq X$ abierto y $u \in \O$, queremos $U_n$ tal que $u \in U_n \subseteq \O$.

    Como el espacio es normal, existe $A \subseteq X$ abierto tal que
    \[u \in A \subseteq \bar A \subseteq \O.\]
    Por lema de Urysohn, existe $f \in C(X)$ tal que $f(x) = 1$ en $A$ y $f(x) = 0$ en $\O^c$.
    Como $\{g_n\}$ es denso, existe un $g_{n_0}$ tal que
    \[\norm{f-g_{n_0}} < \frac{1}{2}.\]
    Si $x \in A$, entonces
    \[g_{n_0}(x) = g_{n_0}(x) - f(x) + f(x) = 1 + g_{n_0}(x) - f(x) > 1 - \frac12 = \frac12.\]
    $\implies U_{g_{n_0}} \subseteq A \implies X$ es metrizable.

    \item \boxed{\impliedby} Supongamos que $X$ es metrizable. Tomamos $d$ como la métrica en $X$. Veamos que $X$ es separable. Me rendí también con la demostración.
  \end{itemize}
\end{proof}

\newpage
\subsubsection{Teorema de Arzela--Ascoli}

\begin{definition}
  Sea $(X,d)$ un espacio métrico. $F \subseteq C(X)$ se dice \wea{equicontinua} si
  \[\forall \e>0, \exists \d>0, \forall f \in F \cor{
    d(u,v)<\d \implies \abs{f(u) - f(v)} < \e
  }.\]
\end{definition}

\begin{definition}
  Decimos que $F \subseteq C(X)$ es \wea{uniformemente acotada} si
  \[\exists M \in \R, \forall f \in F \cor{\norm{f}_\oo < M}.\]
\end{definition}

\begin{theorem}
  [Arzela--Ascoli]
  Si $(X,d)$ es un espacio métrico compacto y $\{f_n\}_{n\in\N} \subseteq C(X)$ es equicontinua y uniformemente acotada, entonces $\{f_n\}_{n\in\N}$ tiene una subsucesión convergente.
\end{theorem}
\begin{proof}
  Como $(X,d)$ es separable, tomamos $D = \{x_n\}_{n\in\N}$ denso en $X$. Seguimos los siguientes pasos:
  \begin{enumerate}[(1)]
    \item Ver que podemos extraer una subsucesión que converge en $D$.
    \item Usar la equicontinuidad para ver que converge en todo $X$.
    \item Ver que el límite obtenido es una función continua.
  \end{enumerate}
  Históricamente, los primeros dos pasos fueron demostrados por Arzelà, y el 3er paso fué un teorema aparte de Ascoli. Actualmente consideramos todo como un único teorema.
  Demostremos cada uno:
  \begin{enumerate}[(1)]
    \item Usaremos un argumento diagonal. Como $\{f_n(x)\}$ es uniformemente acotada en $\R$, tiene una subsucesión convergente $\{f_{s(n,1)}\}_{n \in \N}$. Luego $\{f_{S(n,1)}(x_2)\}$ es uniformemente acotada y tiene subsucesión convergente $\{f_{S(n,2)}(x_2)\}_{n\in\N}$. Notemos que
    \[\{f_{S(n,2)}(x_1)\}_{n\in\N} \subseteq \{f_{S(n,1)}(x_1)\}_{n\in\N}.\]
    Me perdí.
  \end{enumerate}
\end{proof}

\begin{exercise}
  Probar que $(C(X), \norm{\ \cdot \ }_\oo)$ es completo.
\end{exercise}

% \subsubsection{Teorema del Punto Fijo de Banach}
% \subsubsection{Teorema de Categoría de Baire}

%\fecha{Clase 9.}
%\fecha{Clase 10.}
%\fecha{Clase 11.}
%\fecha{Clase 12.}
%\fecha{Clase 13.}
%\fecha{Clase 14.}
%\fecha{Clase 15.}
%\fecha{Clase 16.}
%\fecha{Clase 17.}
%\fecha{Clase 18.}
%\fecha{Clase 19.}
%\fecha{Clase 20.}
%\fecha{Clase 21.}
%\fecha{Clase 22.}
%\fecha{Clase 23.}
%\fecha{Clase 24.}
%\fecha{Clase 25.}
%\fecha{Clase 26.}
%\fecha{Clase 27.}
%\fecha{Clase 28.}
%\fecha{Clase 29.}
%\fecha{Clase 30.}

\end{document}